\documentclass[11pt,a4paper,twoside]{tesis}
% SI NO PENSAS IMPRIMIRLO EN FORMATO LIBRO PODES USAR
%\documentclass[11pt,a4paper]{tesis}

\usepackage{graphicx}
\usepackage[utf8]{inputenc}
\usepackage[spanish]{babel}
\usepackage[left=3cm,right=3cm,bottom=3.5cm,top=3.5cm]{geometry}
\usepackage{listings}
\usepackage{xcolor}
\usepackage{color}
\usepackage{textcomp}

\definecolor{source}{gray}{0.85}% my comment style
\newcommand{\myCommentStyle}[1]{{\footnotesize\sffamily\color{gray!100!white} #1}}
%\newcommand{\myCommentStyle}[1]{{\footnotesize\sffamily\color{black!100!white} #1}}

% my string style
\newcommand{\myStringStyle}[1]{{\footnotesize\sffamily\color{violet!100!black} #1}}
%\newcommand{\myStringStyle}[1]{{\footnotesize\sffamily\color{black!100!black} #1}}

% my symbol style
\newcommand{\mySymbolStyle}[1]{{\footnotesize\sffamily\color{violet!100!black} #1}}
%\newcommand{\mySymbolStyle}[1]{{\footnotesize\sffamily\color{black!100!black} #1}}

% my keyword style
\newcommand{\myKeywordStyle}[1]{{\footnotesize\sffamily\color{green!70!black} #1}}
%\newcommand{\myKeywordStyle}[1]{{\footnotesize\sffamily\color{black!70!black} #1}}

% my global style
\newcommand{\myGlobalStyle}[1]{{\footnotesize\sffamily\color{blue!100!black} #1}}
%\newcommand{\myGlobalStyle}[1]{{\footnotesize\sffamily\color{black!100!black} #1}}

% my number style
\newcommand{\myNumberStyle}[1]{{\footnotesize\sffamily\color{brown!100!black} #1}}
%\newcommand{\myNumberStyle}[1]{{\footnotesize\sffamily\color{black!100!black} #1}}

\lstset{
language={},
% characters
tabsize=3,
escapechar={!},
keepspaces=true,
breaklines=true,
alsoletter={\#},
literate={\$}{{{\$}}}1,
breakautoindent=true,
columns=fullflexible,
showstringspaces=false,
% background
frame=single,
aboveskip=1em, % automatic space before
framerule=0pt,
basicstyle=\footnotesize\sffamily\color{black},
keywordstyle=\myKeywordStyle,% keyword style
commentstyle=\myCommentStyle,% comment style
frame=single,%
backgroundcolor=\color{source},
% numbering
stepnumber=1,
numbersep=10pt,
numberstyle=\tiny,
numberfirstline=true,
% caption
captionpos=b,
% formatting (html)
moredelim=[is][\bfseries]{<b>}{</b>},
moredelim=[is][\textit]{<i>}{</i>},
moredelim=[is][\underbar]{<u>}{</u>},
moredelim=[is][\color{red}\uwave]{<wave>}{</wave>},
moredelim=[is][\color{red}\sout]{<del>}{</del>},
moredelim=[is][\color{blue}\underbar]{<ins>}{</ins>},
% smalltalk stuff
morecomment=[s][\myCommentStyle]{"}{"},
%    morecomment=[s][\myvs]{|}{|},
morestring=[b][\myStringStyle]',
moredelim=[is][]{<sel>}{</sel>},
moredelim=[is][]{<rcv>}{</rcv>},
moredelim=[is][\itshape]{<symb>}{</symb>},
moredelim=[is][\scshape]{<class>}{</class>},
morekeywords={true,false,nil,self,super,thisContext},
identifierstyle=\idstyle,
}

\makeatletter
\newcommand*\idstyle[1]{%
\expandafter\id@style\the\lst@token{#1}\relax%
}
\def\id@style#1#2\relax{%
\ifnum\pdfstrcmp{#1}{\#}=0%
% this is a symbol
\mySymbolStyle{\the\lst@token}%
\else%
\edef\tempa{\uccode`#1}%
\edef\tempb{`#1}%
\ifnum\tempa=\tempb%
% this is a global
\myGlobalStyle{\the\lst@token}%
\else%
\the\lst@token%
\fi%
\fi%
}
\makeatother


%\newcommand{\ct}{\lstinline[backgroundcolor=\color{white}]}
%\newcommand{\needlines}[1]{\Needspace{#1\baselineskip}}
\newcommand{\lct}{\texttt}

\lstnewenvironment{code}{%
    \lstset{%
    % frame=lines,
    frame=single,
    framerule=0pt,
    mathescape=false
    }%
    \noindent%
    \minipage{\linewidth}%
}{%
    \endminipage%
}%


\lstnewenvironment{codeWithLineNumbers}{%
    \lstset{%
    % frame=lines,
    frame=single,
    framerule=0pt,
    mathescape=false,
    numbers=left
    }%
    \noindent%
    \minipage{\linewidth}%
}{%
    \endminipage%
}%



\newenvironment{codeNonSmalltalk}
{\begin{alltt}\sffamily}
{\end{alltt}\normalsize}


\begin{document}


%%%% CARATULA

\def\autor{Fernando Gasperi Jabalera}
\def\tituloTesis{Refactorings automáticos de alto nivel}
\def\runtitulo{Refactorings automáticos de alto nivel}
\def\runtitle{High level automated refactorings}
\def\director{Hernán Wilkinson}
\def\lugar{Buenos Aires, 2020}
\newcommand{\HRule}{\rule{\linewidth}{0.2mm}}
%
\thispagestyle{empty}

\begin{center}\leavevmode

\vspace{-2cm}

\begin{tabular}{l}
\includegraphics[width=2.6cm]{logofcen.pdf}
\end{tabular}


{\large \sc Universidad de Buenos Aires

Facultad de Ciencias Exactas y Naturales

Departamento de Computaci\'on}

\vspace{6.0cm}

%\vspace{3.0cm}
%{
%\Large \color{red}
%\begin{tabular}{|p{2cm}cp{2cm}|}
%\hline
%& Pre-Final Version: \today &\\
%\hline
%\end{tabular}
%}
%\vspace{2.5cm}

\begin{huge}
\textbf{\tituloTesis}
\end{huge}

\vspace{2cm}

{\large Tesis de Licenciatura en Ciencias de la Computaci\'on}

\vspace{2cm}

{\Large \autor}

\end{center}

\vfill

{\large

{Director: \director}

\vspace{.2cm}

{Codirector: \codirector}

\vspace{.2cm}

\lugar
}

\newpage\thispagestyle{empty}


%%%% ABSTRACTS, AGRADECIMIENTOS Y DEDICATORIA
\frontmatter
%%  \pagestyle{empty}
%%  %\begin{center}
%\large \bf \runtitulo
%\end{center}
%\vspace{1cm}
\chapter*{\runtitulo}

\noindent La princesa Leia, líder del movimiento rebelde que desea reinstaurar la República en la galaxia en los tiempos ominosos del Imperio, es capturada por las malévolas Fuerzas Imperiales, capitaneadas por el implacable Darth Vader. El intrépido Luke Skywalker, ayudado por Han Solo, capitán de la nave espacial ``El Halcón Milenario'', y los androides, R2D2 y C3PO, serán los encargados de luchar contra el enemigo y rescatar a la princesa para volver a instaurar la justicia en el seno de la Galaxia (aprox. 200 palabras).

\bigskip

\noindent\textbf{Palabras claves:} Guerra, Rebelión, Wookie, Jedi, Fuerza, Imperio (no menos de 5).
%%  
%%  \cleardoublepage
%%  %\begin{center}
%\large \bf \runtitle
%\end{center}
%\vspace{1cm}
\chapter*{\runtitle}

\noindent In a galaxy far, far away, a psychopathic emperor and his most trusted servant -- a former Jedi Knight known as Darth Vader -- are ruling a universe with fear. They have built a horrifying weapon known as the Death Star, a giant battle station capable of annihilating a world in less than a second. When the Death Star's master plans are captured by the fledgling Rebel Alliance, Vader starts a pursuit of the ship carrying them. A young dissident Senator, Leia Organa, is aboard the ship \& puts the plans into a maintenance robot named R2-D2. Although she is captured, the Death Star plans cannot be found, as R2 \& his companion, a tall robot named C-3PO, have escaped to the desert world of Tatooine below. Through a series of mishaps, the robots end up in the hands of a farm boy named Luke Skywalker, who lives with his Uncle Owen \& Aunt Beru. Owen \& Beru are viciously murdered by the Empire's stormtroopers who are trying to recover the plans, and Luke \& the robots meet with former Jedi Knight Obi-Wan Kenobi to try to return the plans to Leia Organa's home, Alderaan. After contracting a pilot named Han Solo \& his Wookiee companion Chewbacca, they escape an Imperial blockade. But when they reach Alderaan's coordinates, they find it destroyed - by the Death Star. They soon find themselves caught in a tractor beam \& pulled into the Death Star. Although they rescue Leia Organa from the Death Star after a series of narrow escapes, Kenobi becomes one with the Force after being killed by his former pupil - Darth Vader. They reach the Alliance's base on Yavin's fourth moon, but the Imperials are in hot pursuit with the Death Star, and plan to annihilate the Rebel base. The Rebels must quickly find a way to eliminate the Death Star before it destroys them as it did Alderaan (aprox. 200 palabras).

\bigskip

\noindent\textbf{Keywords:} War, Rebellion, Wookie, Jedi, The Force, Empire (no menos de 5). % OPCIONAL: comentar si no se quiere
%%  
%%  \cleardoublepage
%%  \chapter*{Agradecimientos}

\noindent Lorem ipsum dolor sit amet, consectetur adipiscing elit. Fusce sapien ipsum, aliquet eget convallis at, adipiscing non odio. Donec porttitor tincidunt cursus. In tellus dui, varius sed scelerisque faucibus, sagittis non magna. Vestibulum ante ipsum primis in faucibus orci luctus et ultrices posuere cubilia Curae; Mauris et luctus justo. Class aptent taciti sociosqu ad litora torquent per conubia nostra, per inceptos himenaeos. Mauris sit amet purus massa, sed sodales justo. Mauris id mi sed orci porttitor dictum. Donec vitae mi non leo consectetur tempus vel et sapien. Curabitur enim quam, sollicitudin id iaculis id, congue euismod diam. Sed in eros nec urna lacinia porttitor ut vitae nulla. Ut mattis, erat et laoreet feugiat, lacus urna hendrerit nisi, at tincidunt dui justo at felis. Class aptent taciti sociosqu ad litora torquent per conubia nostra, per inceptos himenaeos. Ut iaculis euismod magna et consequat. Mauris eu augue in ipsum elementum dictum. Sed accumsan, velit vel vehicula dignissim, nibh tellus consequat metus, vel fringilla neque dolor in dolor. Aliquam ac justo ut lectus iaculis pharetra vitae sed turpis. Aliquam pulvinar lorem vel ipsum auctor et hendrerit nisl molestie. Donec id felis nec ante placerat vehicula. Sed lacus risus, aliquet vel facilisis eu, placerat vitae augue.
 % OPCIONAL: comentar si no se quiere
%%  
%%  \cleardoublepage
%%  \hfill \textit{A mi persona favorita.}
  % OPCIONAL: comentar si no se quiere
%%  
%%  \cleardoublepage
%%  \tableofcontents
%%  
%%  \mainmatter
%%  \pagestyle{headings}
%%  
%%  
%%  
%%  %%%% ACA VA EL CONTENIDO DE LA TESIS
%%  
%%  \chapter{Introducción}
%%  Resumen general de la tesis. por qué es una pregunta interesante, qué pregunta y respuestas.
%%  
%%  % \section{Infancia y juventud}
%%  % {\begin{small}%
%%  % \begin{flushright}%
%%  % \it
%%  % There's nothing for me now.
%%  % I want to learn the ways of\\ the Force and become a Jedi like my father. \\
%%  % --Luke Skywalker
%%  % \end{flushright}%
%%  % \end{small}%
%%  % \vspace{.5cm}}
%%  
%%  \chapter{Background}
%%  Información necesaria que puedan necesitar los lectores para entender la tesis.
%%  - smalltalk es un entorno vivo de objetos que es reflexivo
- todo es un objeto, las clases son objetos
- cómo funciona el method lookup, cuál es la diferencia entre self y super
- qué son las variables de clase, qué son las variables de pool, qué son las variables de instancia
- qué es una metaclase
- qué es null y cómo se representa en Smalltalk
- qué es un tipo? Es un conjunto de mensajes. Todos los objetos que responden ese conjunto de mensajes
comparten el tipo determinado por ese conjunto de mensajes.
%%  
%%  
%%  \chapter{Estado del arte}
%%  Presentar toda la investigación relevante con respecto al problema que estás investigando.
%%  \section{Introducción a los refactorings, su historia y razón de ser}
% Los proyectos de software evolucionan
% Qué quiere decir que un proyecto evolucione?
% La complejidad de los proyectos de software se incrementa debido a que evolucionan. 
Los sistemas de software evolucionan a lo largo del tiempo. Los requerimientos funcionales y no
funcionales del sistema cambian, requieren mantenimiento e incluso el equipo que lo mantiene cambia.
% REF agregar cita sobre Lehman
La evolución de acuerdo a las leyes informales de Lehman lleva a un incremento de la complejidad del
sistema, en gran parte debido a la erosión de su diseño, lo cual decrementa su mantenibilidad.
% ¿Por qué la evolución incrementa la complejidad?
La erosión del diseño se produce porque no es posible anticipar los cambios que se le realizarán al
software en el futuro entonces el diseño original no es apropiado para incorporar todos los cambios
que se le realizan luego al sistema. El desarrollador no modifica el diseño antes de extender o
modificar el modelo, entonces las modificaciones que le realiza lo vuelven más complejo. Los
proyectos de sotware industriales tienen restricciones de tiempo ajustadas entonces los desarrolladores
introducen modificaciones de la forma más veloz posible sin tener en cuenta la pérdida de calidad.
% WRITE El incremento de la complejidad reduce su mantenibilidad. ¿Qué consecuencias negativas trae esto?
% ¿Qué técnicas existen para lidiar con este problema?
Existen numerosas técnicas
para minimizar o reducir el impacto negativo del cambio en los sistemas de software. Una metodología
general para desarrollar software que se basa en la aceptación del cambio del mismo es la conocida
% REF agregar referencia a extreme programming
extreme programming, la cual utiliza técnicas como:

\begin{itemize}
    \item scrum
    \item tdd
    \item unit testing
    \item refactoring
    \item pair programming
    \item sprints
    \item etc.
\end{itemize}

% ¿Qué son los refactorings?
El refatoring del software es una de ellas. Se utiliza en técnicas de desarrollo iterativo
incremental, no se han visto integradas en sistemas tradicionales con modelos de desarrollo de cascada lineales.
En el contexto de las metodologías ágiles los refactorings se realizan durante el desarrollo del
software como una parte integral del mismo y no como una etapa separada. El término fue acuñado por
% REF agregar referencia a la tesis de Opdyke
Opdyke en su tesis de doctorado y luego popularizado por Fowler, uno de los mayores evangelistas de
% REF agregar referencia a el libro de Fowler
la metodología extreme programming, en su libro \textit{Refactoring: Improving the design of existing code}.
Un refactoring o reestructuración es una modificación al software que no cambia su funcionalidad. El objetivo
de la misma es mejorar la calidad del sistema modificando su estructura interna para volverlo más mantenible,
entendible o que se adapte mejor a futuros cambios o funcionalidades que haya que agregarle. Por
ejemplo, si se desea agregar una funcionalidad y el diseño no la contempló originalmente se puede
agregar al diseño actual o cambiar el diseño primero para que agregar la funcionalidad sea más
simple y se pueda seguir extendiendo el software en esa dirección más fácilmente.


\section{Investigaciones sobre software refactorings}
Los software refactorings son investigados desde numerosas perspectivas. Alguno de los temas más
investigados son:
% REF a Trends, opportunities...
\begin{itemize}
    % REF a el paper que definía floss refactoring
    \item refactoring tools: qué factores afectan el uso de refactorings automáticos? qué
refactorings automáticos usan los programadores? cuáles no y por qué? cómo mejorar las herramientas actuales?
    \item la relación entre los bad smells y los refactorings que podrían ayudar a lidiar con ellos. No hay estudios
empíricos que lleguen a conclusiones claras sobre la utilidad de los code smells para lo
programadores o si ayudan a determinar que un refactoring es necesario en ese lugar o qué
refactoring es necesario.
    \item refactorings en artifacts que no son código: especificación de requerimientos, diseño de más alto nivel, documentación, etc.
    \item patrones de diseño: su utilidad, cuánto ayudarían, cuán de alto nivel se pueden hacer es una interrogante
    \item TDD, refactoring es una importante etapa de esta técnica 
    \item métricas del software sobre los atributos de calidad para medir el impacto de los
refactorings  
\end{itemize}


\section{Complejidades de estudiar refactorings}

\begin{itemize}
    \item cómo comparar las herramientas y los refactorings? taxonomía
    \item cómo determinar si un refactoring preserva el comportamiento? poco formal, la mayoría de
    los enfoques son semi formales o solo con los tests del sistema
    \item entender cómo realizan refactorings los desarrolladores actualmente? difícil, los más
    precisos son poco representativos, los más generales son poco precisos y solo uno analiza los
    cambios en continuo.
    \item métricas para determinar si un refactoring aumentó la calidad del software o no? redujo la
    complejidad? tampoco está claro cómo hacerlo.
\end{itemize}

\section{¿Para qué realizar cambios que no impactan en la funcionalidad del software?}
El software evoluciona porque a
medida que el tiempo pasa los requerimientos que intenta satisfacer cambian. Para satisfacer los
nuevos requerimientos es necesario modificar y extender el sistema, además de arreglar los posibles
errores que surjan. Los atributos de calidad más relacionados con esta parte del desarrollo del
software son: 1) extensibilidad 2) mantenibilidad 3) comprensiblidad.

\section{¿Son útiles los refactorings?}
Los refactorings son útiles porque permiten preservar la calidad del software, principalment su: 1)
extensibilidad 2) su mantenibilidad y 3) su comprensibilidad.
(Practical Analysis for Refactoring)
Extreme Programming acepta que el software evoluciona y el diseño constantemente cambia. La
adaptación del software requiere energía del desarrollador constante y es necesaria para que la
complejidad no crezca demasiado. Cambiar el diseño aplicando refactorings para que la complejidad
crezca lo menos posible solo será posible minimizando la energía que tiene que poner el
desarrollador y esto se puede conseguir automatizando los refactorings.
(Empirical study on the impact...)
Los refactorings aumentan la mantenibilidad del software. ¿Qué relación tienen con la
calidad del software? Permiten preservarla. Los 3 atributos mencionados en el punto anterior se van
degradando a medida que el software evoluciona y se produce un fenómeno conocido como erosión del
diseño. Los cambios que se introducen subsecuentemente en el software no fueron previstos por el
diseñador original del sistema y en consecuencia la calidad decrece. Los refactorings modifican el
diseño antes de agregar las nuevas funcionalidades para que el diseño pueda adoptar los nuevos
cambios sin perder calidad. ¿Son útiles los refactorings automáticos? ¿qué se sabe? Detectar cuándo
realizar un refactoring, dónde realizarlo y qué cambiar son todos problemas íntimamente
relacionados. Existen investigaciones que intentan procesar el código y realizarle cambios sin
intervención del desarrollador. Sin embargo, la mayoría de las herramientas automatizan la ejecución
del cambio, el desarrollador elige cuándo, dónde y qué cambio se realizará. La cantidad de
refactorings posibles es grande por lo cual se estima que sería útil contar con una herramienta que
le permita al desarrollador crear sus propios refactorings automatizados y luego utilizarlos
(Trends...). Hay pocos estudios en referencia a code smells (Trends...), la detección de los mismos
y su utilización para elegir un refactoring no tiene atención.
(A Refactoring Tool for Smalltalk)
Las modificaciones que se le realizan a un programa pueden ser divididas en dos etapas:
1) refactoring del código/diseño para adaptarse mejor a las modificaciones
2) las modificaciones. Si se cuenta con refactorings automáticos que preservan el comportamiento
para realizar la primer etapa entonces la única fuente de errores al introducir cambios en un
programa es la segunda etapa. Esto reduce la posible cantidad de errores agilizando el mantenimiento
del software.
(Programmer-Friendly...) Los dos principales riesgos al aplicar refactorings son:

- introducir bugs - tiempo consumido

los refactorings automáticos mitigan ambas.

(Automated refactoring of super-class method invocations to the Template Method design pattern) Los
principales caminos que se están investigando actualmente para la automatización de la
identificación de diseños pobres y sus correspondientes refactorings sin intervención del
desarrolladors son:

\begin{itemize}
    \item metric-based: low quality parts of code are identified by detecting the minimums of some
    quality metric.
    \item logic-based: the code is translated to an intermediate logic language which is analyzed
    with logic rules to identify flaws.
    \item search-based: the improvement of the design is translated into an optimization problem of
    a fitness function whose search space consists of the alternative designs.
    \item visualization techniques: differents forms of visualizing the code aim to aid the
    developer to gain new perspectives of the code that would allow him to spot flaws more easily.
\end{itemize}

Categorización de los métodos utilizados para realizar refactorings automáticos

\begin{itemize}
    \item método de aplicación: ¿cuán automático es? automatiza identificación, elección de nombres, aplicación del refactoring, cuándo, dónde.
    \item preservación del comportamiento: manual, semi-formal, formal.
    \item composición de los refactorings: dinámica o estática.
\end{itemize}

\section{¿Son útiles los refactorings de alto nivel? ¿qué se sabe?}
Existen opiniones encontradas. Algunos
opinan que cuanto más se automatice mejor ya que modificar el código es una operación delicada,
menor intervención humana implicaría menor espacio para el error. Sin embargo, estudios acerca
de la utilización de estas herramientas no arrojan resultados claros que indiquen que los
desarrolladores las utilicen con la frecuencia esperada. Incluso la correlación entre la
complejidad del refactoring (cuán alto es su nivel) y la frecuencia de su utilización se muestra
inversa. No queda claro a qué se debe esto. Hay estudios que intentan entender a qué se debe
esto y cómo construir herramientas que automaticen los cambios y que los desarrolladores
utilicen. (A discussion of software..) Los refactorings deben ser más complejos para poder
ayudar al desarrollo de proyectos grandes. Por lo tanto, las herramientas deberían permitir
componer refactorings para poder construir más complejos. (Practical Analysis for refactoring)
La aplicación de patrones de diseño se correspondería con refactorings compuestos que
permitirían ahorrar mucha energía del desarrollador y explorar opciones de diseño con un costo
menor. (A Refactoring Tool... A Compositional Paradigm) Los refactorings complejos deberían
realizarse como la composición de refactorings más simples automáticos. Algunos refactorings se
realizan comúnmente y es tedioso realizarlos incluso con refactorings automáticos más simples
porque no es fácil encontrar refactorings intermedios que simplifiquen la tarea (por ejemplo,
introduce null object, cómo lo partirías en refactorings simples?) ¿Qué forma toman los
refactorings de alto nivel? por qué? Los refactorings de alto nivel son combinaciones comunes de
refactorings más simples que pueden o no tener una semántica clara a nivel diseño del sistema.
Los patrones de diseño son soluciones a problemas de diseño que surgen frecuentemente al
construir sistemas con lenguajes orientados a objetos. Los patrones de diseño son soluciones
vagamente definidas al nivel del diseño de un sistema que le permiten soportar cierta
funcionalidad con más calidad. Su presencia es común y es por eso que los refactorings de alto
nivel automáticos intentan aplicar estos patrones generalmente. Hay varios ejemplos de papers.

(A Methodology for the Automated Introduction of Design Pattern) Los refactorings de alto nivel
tienen más puntos de partida posibles y más destinos posibles ya que, a diferencia de refactorings
más primitivos, están definidos de forma más relajada. Se puede decir que cuanto más alto el nivel
de un refactoring menos específica es su definición. Se han clasificado los distintos puntos de
partida posibles: 

- hoja en blanco: las entidades que se relacionarían en el patrón de diseño no se conocen todavía.
Este caso no ocurre en la práctica usualmente. - anti-patrón: este caso se debe a falta de
conocimiento del programador. Se soluciona con educación, los posibles malos diseños son demasiados
para considerarlos uno por uno. - precursor: es un buen diseño para un caso más simple pero que ante
nuevas necesidades de extensión debe cambiarse.

\section{¿Por qué refactorings automáticos que introducen patrones de diseño?}
Al automatizar refactorings se
quiere automatizar lo más posible para ahorrar la mayor cantidad de energía y tiempo del
desarrollador, y al nivel más expresivo posible para que la aplicación del refactoring esté lo más
cerca del nivel de abstracción al que está pensando el desarrollador el cambio. Las fuerzas que se
contraponen son expresividad y precisión. Cuanto más alto el nivel de abstracción, más difícil es
precisar a nivel de código en qué consiste el cambio. Los patrones de diseño presentan un balance
atractivo porque son cambios semántico suficientemente específicos para precisarlos en el código que
además están cerca de la forma que tiene el desarrollador de pensar su cambio.

¿En qué escenarios se utilizarían estos refactorings? El desarrollador tiene que agregar una nueva
funcionalidad al sistema pero su diseño actual no la acomoda de manera simple. El desarrollador
identifica dónde puede realizar un cambio para lograr la flexibilidad deseada y qué cambio quiere
realizar (la introducción de un patrónde diseño). En ese momento se puede automatizar la aplicación
del refactoring. patrones de diseño (null object pattern paper) Se ha explorado la introducción
automática de patrones de diseño orientada a root canal refactoring, ya que son herramientas que en
batch analizan todo el código, presentan los candidatos identificados y proveen la opción de aplicar
el refactoring. Las investigaciones se pueden clasificar según el tipo de patrones que analizaron:

- estructurales (Abstract Factory y Composite) - de comportamiento (Decorator, Template Method, Null
Object y State/Strategy)

los métodos que utilizan para la identificación de oportunidades de mejora al diseño y la aplicación
de los respectivos refactorings también es variada. Sin embargo, el principal problema que
le vemos a estos trabajos es que no se ajustan a la forma de trabajo del programador.
refactorings compuestos

\section{Preservación del comportamiento}
Este punto está presente en la misma
definición de un refactoring y es de máxima importancia. Por el momento la preservación de
comportamiento no se puede definir de una forma que permita implementarla en las
herramientas. La preservación de la funcionalidad es difícil de formalizar y existen
distintos tipos de software para los cuales preservar el comportamiento implica más que
preservar su funcionalidad, por ejemplo:

- tiempo de ejecución (sistemas de tiempo real) - memoria utilizada y consumo de energía (sistemas
embebidos) - condiciones de seguridad (sistemas en los cuales la seguridad es crítica)

Por esta razón un testing sistematizado y ajustado a los requerimientos de cada sistema particular
es la única y mejor herramienta con la que se cuenta actualmente.

(A Refactoring Tool for Smalltalk) No es posible asegurar que se preserva el comportamiento en un
lenguaje con tipado dinámico como Smalltalk, la única forma de converger a un programa correcto es a
través del análisis dinámico. Esto no lo realicé pero incluso el análisis dinámico se basa en una
suite de tests representativa. (Improving refactoring tools in Smalltalk using static type
inference) La única forma de asegurar que los refactorings son correctos es con pruebas formales.
Las herramientas actuales no realizan esto porque la complejidad del software actual vuelve
demasiado costoso aplicar modelos de verificación formal a los programas que nos interesa
reestructurar.

(Automated Application of Design Patterns: A Refactoring Approach) Perspectivas sobre la
preservación de comportamiento Las principales técnicas que se utilizan para lidiar con la
preservación de comportamiento son:

- no prestarle atención - semi formal para tener una referencia de lo pensado y confianza en ello -
verificación formal

No se utiliza la verificación formal porque pocos lenguajes de programación ampliamente utlizados
    tienen una semántica formal y un compilador que la verifica. Además, la complejidad de las
    demostraciones de preservación de comportamiento para transformaciones no triviales es
    intratable. Refactoring implementados en IDEs populares Enumerar y compara las implementaciones
    de refactorings en IDEs. Eclipse IntelliJ
%    [Explicación](https://www.jetbrains.com/help/idea/extract-into-class-refactorings.html#extract_method_object)
    superficial de IntelliJ Pasos Elige el código Elige Delegate, Parameter, Method Elige nombre de
    la clase y del método Elige opciones sobre los parámetros Prevee los cambios y los confirma No
    se puede elegir cualquier parte del código, tiene que ser un conjunto de statements o
    expressions No parametriza el contexto de la clase, solo lo local NetBeans no tiene
    ExtractToMethodObject tiene extract method tradicional Visual Studio CodeRush Visual Assist Wing
    IDE XCode Extract to method object IntelliJ Idea Solo extrae a una clase interna, lo cual
    simplifica bastante el refactoring porque compartimos el scope de variables de instancia, de
    clase, etc. Visual Studio Solo realiza extract method y el nuevo método pertenece a la misma
    clase. Si necesita devolver más de un resultado utiliza los parámetros de salida de C\#. Eclipse
    Solo realiza extract method.
%%  
%%  \chapter{Objetivos y justificación}
%%  \section{Justificación}
\subsection{¿Por qué refactorings automáticos de alto nivel?}
% Preguntas guía
%   - Cuál es el potencial de los refactorings automáticos de alto nivel
%   - Por qué creemos que son una buena herramienta

% re redactar este párrafo, las ideas están bien pero no fluye. Cada oración
% empieza repitiendo el sujeto anterior, parece que lo escribió un robot.
Los refactorings son una herramienta esencial del arsenal de los ingenieros de
software para minimzar el incremento de la complejidad que implica la evolución
del software y así preservar su mantenibilidad. Sin embargo, los
beneficios de los refactorings conllevan riesgos y costos:

\begin{description}
    \item[Introducir bugs] Los refactorings automáticos permiten reducir el
      riesgo de introducir bugs porque las herramientas que los aplican
      recibieron más tiempo de desarrollo, testeo y consideración que la que
      cualquier desarrollador puede dedicar a la hora de realizar un refactoring
      manual. Además, cuentan con varias ventajas sobre la aplicación manual:
        \begin{itemize}
            \item cuanto más se usan mejores se vuelven: los errores que encuentren los
            usuarios se pueden arreglar para converger a través del tiempo a una
            preservación de comportamiento con una probabilidad altísima. Los
            programadores pueden cometer el mismo error repetidas veces; algo que
            contribuye a esto es que la experiencia de los programadores trabajando en
            el proyecto puede variar y el tiempo que hayan pasado trabajando en ese
            proyecto en particular también.
            \item no son afectadas por el contexto: los refactorings manuales son
            realizados por desarrolladores a los cuales los puede afectar el
            cansancio, el estrés, el apuro, etc. factores que aumentan la probabilidad
            de equivocación u olvido en un tarea tediosa y demandante de atención al
            detalle como esta.
        \end{itemize}
    \item[Tiempo de desarrollo] El tiempo de desarrollo se vería disminuido
      porque una herramienta automática es más veloz que un humano, si esto no
      fuera así la herramienta es inútil. Esto representa una gran ganancia ya
      que el mantenimiento de un sistema es su mayor costo, si este costo se
      disminuye se abarataría la producción de software, y, considerando la
      omnipresencia actual del mismo, la herramienta tendría un impacto enorme
      en la sociedad.
    \item[Reducir la calidad del sistema] Al no contar con métricas definitivas
      para evaluar la calidad de un sistema no es simple ver si efectivamente
      los refactorings automáticos mitigarían el riesgo de bajar la calidad del
      sistema. Sin embargo, existen consideraciones que nos llevan a intuir que
      este puede ser el caso.
      % REF al que decía que con refactorings de patrones iban a poder probar más diseños
      UNA REF postula que si los desarrolladores contaran con buenas herramientas
      automáticas para realizar refactorings de alto nivel probarían más opciones de
      diseño, lo cual les permitiría llegar a mejores diseños. No es concluyente pero
      si ciertas herramientas aportan a que los desarrolladores lleguen a mejores
      diseños tienen una buena chance de aumentar la calidad de los sistemas.
\end{itemize}


Los refactorings automáticos de alto nivel cuentan con el potencial de maximizar
estos beneficios porque el scope de automatización sería mayor, lo cual implica
más ahorro de tiempo de desarrollo y menos posibilidad de introducir bugs porque
existe más area del refactorings cubierta, y porque al operar a un nivel de
abstracción mayor, más cercano al nivel de la semántica del pensamiento del
programador, aumentaría la velocidad a la que él probar diseños alternativos y
así llegar a uno mejor.

% Por qué refactorings automáticos de alto nivel
% REF hay un empirical study de esto
[La referencia] concluye que los refactorings automátricos de alto nivel son más
escalables y performantes.  También que los refactorings automáticos aumentan la
velocidad de desarrollo, resultando en una reducción de costos y un aumento de
la calidad del software porque los desarrolladores tienen la posibilidad de
probar más cambios con menos esfuerzo.

Al explorar refactorings de alto nivel surge la pregunta de cuán alto es
convienente subir.
% Los dos objetivos contrapuestos de los refactorings
% REF (Refactoring Object-Oriented Frameworks)
(Refactoring Object-Oriented Frameworks) expone dos objetivos de los
refactorings automáticos que se contraponen. Un objetivo es que el refactoring
sea expresivo, que su nivel sea lo más alto posible para mantener al usuario lo
más lejos posible de toda la complejidad que implica realizar el cambio, y el
otro es que preserve el comportamiento. Cuanto más alto es su nivel más difícil
es justificar que preservan el comportamiento del programa, por lo tanto
queremos que los refactorings sean suficientemente pequeños para poder tener una
buena confianza de que preservan el comportamiento y suficientemente abstractos
para que sean útiles.

% Preguntas que se hace este trabajo, por qué es exploratorio
Este trabajo busca sentar las bases para poder analizar este tipo de
refactorings.



\subsection{Relación con las investigaciones más cercanas}
% Preguntas guía
%   - cuáles son los trabajos más parecidos? los de patrones
% REF A Methodology for the Automated Introduction of Design Pattern
Este paper es similar porque busca implementar un refactoring que introduce un
patrón de diseño. El foco está puesto en la identificación de transformaciones
parciales que permitan reestructurar el programa desde un estado precursor
elegido hasta uno con el patrón aplicado. En parte porque su principal objetivo
es en el futuro poder automatizar la generación de las transformaciones
necesarias. Este objetivo es atractivo pero no necesariamente valioso. Todavía
no contamos la información necesaria para saber si ese nivel de abstracción es
apropiado para los desarrolladores y si esas transformaciones son
generalizables.

% REF
%   - Automated refactoring to the Strategy design pattern
%   - Automated refactoring of super-class method invocations to the Template Method design pattern
%   - Automated refactoring to the Null Object design pattern
Estos 3 papers implementan la introducción de patrones de diseño en
Java. Además, automatizan la identificación de oportunidades para aplicar el
refactoring. La herramienta evalua toda el código, identifica partes del mismo
que se beneficiarían con la aplicación del refactoring y ofrecen aplicarlo. Los
dos problemas principales
% REF a programmer-friendly refactoring tools
que se pueden encontrar son que el nivel de automatización no se ajusta a las
prácticas actuales de desarrollo y que las restricciones impuestas al código
precursor para poder aplicar el refactoring son demasiado estrictas, perdiéndose
así un gran número de casos.

%   - qué no cubren esos trabajos?
%   - cuáles son las diferencias con este trabajo?


\section{Valor de realizar este trabajo}

% Cuáles son los objetivos de este trabajo?
El objetivo de este trabajo es implementar refactorings automáticos de alto
nivel para:

\begin{itemize}
    \item entender mejor sus limitaciones en relación al tradeoff expresividad
      vs precisión.
    \item analizar los desafíos que encontremos para guiar desarrollos futuros.
    \item contar con una herramienta de base que nos permita medir, analizar y
      evaluar cómo se relaciona el programador con la misma y cómo podría
      mejorarse para maximizar su productividad.
\end{itemize}

% Por qué es necesario realizar una exploración de refactorings automáticos de alto nivel?
Se necesitan implementaciones de refactorings de este tipo para poder evaluar
sus potenciales beneficios y así poder mejorar las herramientas que los
aplican. Las IDEs más utilizadas y poderosas no proveen actualmente refactorings
automáticos de este tipo, los disponibles son más simples. Vemos en esto una
oportunidad de mejora y una necesidad de crecimiento grande.

Orientamos esta herramienta a la automatización de la aplicación del refactoring
porque el principal problema que vemos con las herramientas que automatizan más
etapas es que no requieren intervención del programador y por lo tanto no se
ajustan a su forma de trabajo es que corren dos riesgos grandes:

\begin{itemize}
    \item baja frecuencia de uso: nadie los usaría ya que asumen que son parte
      de una etapa del desarrollo que no existe.
    \item incrementar los costos de un proyecto: en primer lugar, en términos de
      tiempo de desarrollo. El problema es que si el refactoring toma decisiones
      buenas pero que perjudican la mantenibilidad del código por hacerlo más
      complejo o menos entendible, como luego el proyecto lo tienen que seguir
      interviniendo programadores humanos cada vez les tomará más tiempo
      realizarle modificaciones.
\end{itemize}

Mientras no haya una manera precisa de expresar comprensibilidad y diseño del
dominio la intervención humana será necesaria y proveerá una reducción de costos
mayor que la utilización de
% REF la survey que decía que el 80% del costo del software es mantenimiento
refactorings completamente automáticos. Aportar en esta dirección aporta valor
porque el mantenimiento del software es su mayor costo.



% Beneficios adicionales de realizar este trabajo:
% % Por qué este trabajo los cumple?
% 
% % Por qué es útil contar con una solución al problema?
% \begin{itemize}
%     \item Chequear precondiciones y postcondiciones de refactorings intermedios es más costoso y
%     complicado. Si esto funciona tendríamos una opción más eficiente.
%     % Algo sobre que no es fácil identificar refactorings intermedios para componerlos y hacer retrofitting de un patrón
%     % WRITE había un paper que propone exactamente esto, el de compositional paradigm, revisar y ponerlo
%     % para que no parezca que lo olvido.
%     \item No está claro como identificar refactorings intermedios que sean útiles y que su
%     composición produzca el retrofitting de un patrón.
%     % REF (A Refactoring Tool for Smalltalk)
%     En (A Refactoring Tool for Smalltalk) se pueden ver ejemplos de retrofitting the patrones de
%     diseño en los que la cantidad de cambios es grandes y no es fácil identificar refactorings de
%     tamaño intermedio que permitan simplificar el cambio. Contar con refactorings sin necesidad
%     de las transformaciones intermedias ahorraría trabajo de desarrollo.
% \end{itemize}


% parte colgada pero que creo que va en esta sección final

% En resumen, queremos analizar cuán útiles son estos refactorings y cuánto contribuyen a la productividad de los
% programadores. Además, queremos analizar la definición de este tipo de refactorings, si se pueden
% precisar sus puntos de partida o si el hecho de que sean tan abstractos los vuelve demasiado
% específicos y ya no cubren una buena superficie de situaciones como para tener un impacto
% significativo.
% Buscamos entender cuáles son los tradeoffs que presentan este tipo de refactorings, qué dificultad
% tiene comunicarle al usuario la acción que está realizando y los alcances de la misma y qué tradeoff
% hay entre configurabilidad y automatización.


%%  \chapter{Descripción de la herramienta}
%%  \section{Decisiones sobre la implementación}

\subsection*{¿Por qué Smalltalk?}
Necesitamos que sea un lenguaje orientado a objetos ya que vamos a estar introduciendo refactorings
que son propios de el diseño orientado a objetos. Un lenguaje dinámico y puro permitirá iterar más
rápidamente sobre la herramienta, por lo tanto es ideal para la realización de un prototipo y nos
permite explorar con más soltura. Al ser un trabajo exploratorio no queríamos vernos limitados por
las particularidades de un lenguaje de programación o que las mayores dificultades provengan de la
navegación de una sintaxis compleja y sobrecargada. Smalltalk cuenta con una sintaxis simple y al
ser un lenguaje de objetos puro y reflexivo provee la posibilidad de manipular los programas a
través de metaprogramación de una manera simple y poderosa. La reflexividad de Smalltalk además (A
Refactoring Tool for Smalltalk) permite utilizar el mismo lenguaje que se está modificando para
modificarlo, no es necesario un metalenguaje intermedio.
% REF (A Refactoring Tool for Smalltalk)
Smalltalk ya ha probado ser una buen campo de experimentación para refactorings (A Refactoring Tool
for Smalltalk)
% REF (A Refactoring Tool for Smalltalk)


\subsection*{¿Por qué CuisUniversity?}
% REF a cuis y a OpenSmalltalk VM
Cuis es un entorno enfocado en Smalltalk-80 gratuito y Open Source. Está basado en la OpenSmalltalk
% REF a cuisUniversity
VM, como Squeak, y su principal objetivo es mantenerse su base pequeña y simple. CuisUniversity es
un ambiente basado en Cuis creado con el objetivo de enseñar programación orientada a objetos.
Actualmente, es utilizado en la Universidad de Buenos Aires, en la Universidad de Quilmes y en la
Universidad Católica Argentina. Uno de los objetivos de este trabajo es que permitir a los alumnos
experimentar con este tipo de refactorings.


\subsection*{¿Por qué Introduce Null Object y Extract Method To Method Object?}
Luego de analizar varios refactorings de alto nivel, la mayoría variaciones de patrones de diseño,
nos enfocamos en estos dos porque creemos que presentan un buen balance de abstracción y
simplicidad. Son lo suficientemente abstractos para ser considerados de alto nivel, no son
refactorings simples que ya estén presentes en otras IDEs ni manipulan elementos diminutos como un
nombre ya que operan al nivel de diseño del código. Sin embargo, son más simples que otros, lo cual
nos permitirá explorar y aprender para poder sentar la bases hacia refactorings futuros. El
Introduce Null Object es un caso particular de Introduce Special Case y de Replace If With
Polymorphism, por lo tanto consideramos que es un buen precursor de ellos. De la misma forma,
Extract Method To Method Object es un precursor de todos los refactorings que requieren extraer
código y hacer reemplazos contextuales.


\section{Los refactorings}
% Límites del scope
% WRITE explicación de lo que describe este párrafo
Varios cambios que formarían parte del refactoring si se realizara manualmente quedaron fuera
del scope porque no vimos fácilmente cómo automatizarlo o por que dadas las limitaciones de
tiempo no entraron dentro de la priorización. Dejamos asentadas acá qué cosas quedaron fuera y
en qué categorías se encuentran. Además, explicamos cómo creemos que podrían resolverse en el
futuro o cómo hacemos nosotros para guiar al usuario en las tareas manuales que debería realizar
luego de aplicar el refactoring.

\chapter{Extact Method to Method Object}
\section{Extract Method to Method Object}

% WRITE Descripción del refactoring
% historia, quién lo inventó?
% qué code smells motivan el introduce null object
un método demasiado largo, con demasiadas responsabilidades. Es poco claro, difícil de testear y de
modificar. No se puede aplicar extractMethod porque hay demasiadas variables temporales
relacionadas.
% en qué consiste?

% WRITE Ejemplo concreto
% WRITE Precursor del refactoring en referencia al paper de Methodology

% Limitaciones
- solo extrae métodos completos
Creemos que agregar un return explícito y parametrizar self en un caso en el que no es necesario
agregaría complejidad innecesaria al refactoring. Por lo tanto, elegimos preservar el return self
implícito del método. Si encontramos que el método a refactorear cuenta con un return self implícito
no incluimos el return de la evaluación del MethodObject para preservar el return self implícito:

% Decisiones tomadas
- el método tiene un return self implícito

\begin{code}

unaryMessageSelector

 1 + 1.

unaryMessageSelector
 
 (MethodObject new) value.

unaryMessageSelector

 ^(MethodObject with: self) value.

\end{code}

% Corte de scope
- agregar un setter de la variable de clase

%%  \chapter{Introduce Null Object}
%%  \section{Motivación y smells}

Existen pocos lenguajes de programación orientados a objetos, y de uso extendido ninguno, que no
cuenten con una u otra forma de representar el valor especial null. La utilización de null conlleva
como principal riesgo el enviarle un mensaje o derefenciarlo, dependiendo el lenguaje utilizado, lo
cual resulta en un error en runtime o una terminación abnormal del programa. La forma más simple de
prevenir estos errores es realizar chequeos sobre las variables que se sospecha pueden contener null
para evitar enviarles mensajes. Veamos un ejemplo en la clase \lstinline{Behavior} de Cuis:

\begin{code}
Behavior>>allSuperclasses
	"Answer an OrderedCollection of the receiver's and the receiver's  
	ancestor's superclasses. The first element is the receiver's immediate  
	superclass, followed by its superclass; the last element is Object."
	| temp |
	^ superclass
		ifNil: [OrderedCollection new]
		ifNotNil: [temp := superclass allSuperclasses.
			temp addFirst: superclass.
            temp]
\end{code}

En este caso se verifica si la variable de instancia \lstinline{superclass} es nil para evitar
enviarle el mensaje \lstinline{allSuperclasses} y además por que la ausencia de superclase indica
que estamos en la raíz de la jerarquía.
% REF a Fowler
Si null conlleva tantos riesgos, ¿por qué se utiliza con tanta frecuencia? La utilización de null
puede ser la forma más simple y barata de implementar un modelo en determinado estado o caso. Sin
embargo, si los chequeos por null abundan las causas más comunes son:

\begin{itemize}
    \item que la variable de instancia sea lo que se llama un temporary field, un colaborador que no
      forma parte de la esencia del objeto y por lo tanto no es necesario siempre; en los casos en los
      que no es necesario es nil.
    \item que falte una abstracción en el modelo para representar la ausencia de un tipo de objetos.
\end{itemize}

% REF a fowler y a woolf
Fowler y Woolf describen que la motivación para realizar el refactoring son repetidos chequeos por
nil de un colaborador. Los condicionales reducen la entendibilidad del código y al chequear siempre
la misma condición producen código repetido, lo cual suele aumentar la posibilidad de que el
programador introduzca bugs olvidándose de chequear por nil antes de enviarle un mensaje al
colaborador correspondiente. Una de las principales herramientas del paradigma de objetos para
lidiar con la proliferación de condicionales repetidos en el código es la utilización de
polimorfismo. En lugar de chequear manualmente en qué tipo de situación nos encontramos para decidir
qué hacer, se envía un mensaje polimórfico a un objeto que realizará la tarea correspondiente para
esa situación dependiendo de su tipo. La introducción de un Null Object sigue esta guía. Se
reemplazará la utilización de nil por un Null Object, un objeto que representa la ausencia o la
nada del tipo correspondiente y encapsula el comportamiento necesario en las situaciones en las
cuales hay una ausencia de ese tipo. Al contar con un Null Object ya no es necesario verificar si el
objeto es nil, nunca debería ser nil, y se le puede enviar los mismos mensajes que a un Concrete
Object, ya que son polimórficos, contando con que el Null Object sabrá qué hacer. Además de reducir 
la complejidad del código la introducción del Null Object consigue que las extensiones futuras sean
más seguras, ya que no hay que chequear por nil, y más mantenibles porque si se requiere cierto
comportamiento adicional en un caso en el que el colaborador está ausente el cambio quedará
encapsulado en la clase del Null Object.


% En el paper de null object hay una buena introducción, si estoy falto de ideas puedo sacar de ahí.


\section{¿En qué consiste?}

La esencia del refactoring es el reemplazo de condicionales que chequean si una variable es nil
o no por el envío de un mensaje polimórfico. Conseguir ese objetivo requiere de varios pasos que tienen
varias alternativas de implementación:

\begin{enumerate}
    \item Crear una jerarquía para el Null Object.
    \item Reemplazar las asignaciones de nil por asignaciones del Null Object.
    \item Reemplazar los condicionales que chequean nil por envíos polimórficos al objeto.
    \item Mover el comportamiento requerido para los casos en los que el objeto es nil al Null
    Object.
\end{enumerate}

Llamaremos NullComponent a la clase del Null Object y Component a la clase del objeto que representa
un componente que no es nil. El NullComponent será polimórfico con el Component pero no contiene
estado porque sus métodos no tienen comportamiento, están vacíos, o devuelven alguna constante, en
ocasiones en forma de un Null Object de otro tipo. En un lenguaje dinámico como Smalltalk no es
necesario que exista una relación de herencia entre las clases para que puedan utilizarse
indistintamente como en el caso de lenguajes estáticamente tipados. De todas formas se elige
utilizar una relación de herencia y relacionar al Component y al NullComponent en una misma
jerarquía para dejar la relación entre las dos explícita. La jerarquía que los relaciona puede tomar
varias formas:

% NullObject hierarchy diagram 

La primer y más simple alternativa es la propuesta por Fowler. La principal ventaja de esta
alternativa es que no requiere cambiar la clase Component; su principal desventaja es que si se le
agregan mensajes a la interfaz de Component el NullComponent los heredará y si no se recuerda
reescribirlos en el NullComponent con el comportamiento correcto para el NullComponent el programa
tendrá comportamientos inesperados. La segunda alternativa no tiene la desventaja que vimos en la
primera y es por eso que suele preferirse por sobre ella.

Luego de armar la jerarquía se deberán sobrescribir todos los mensajes de la interfaz de Component
en NullComponent con un comportamiento vacío. Luego de realizar esto se reemplazarán las asignaciones de
nil a la variable de instancia. Por ejemplo:

\begin{code}
Context>>initialize
    component := null.

"Pasará a ser"

Context>>initialize
    component := NullComponent new.
\end{code}

Algunas asignaciones pueden tener como lado derecho el envío de un mensaje el cual puede devolver 
nil. Esos casos deben reemplazarse por una delegación que permita asignar el Null Object:

\begin{code}
Context>>sendTask
    component := self latestComponent.

"Se reemplazará por:"

Context>>assignComponentWith: aComponentOrNull
    component := aComponentOrNull ifNil: [NullComponent new] ifNotNil: [aComponentOrNull].

Context>>sendTask
    self assignComponentWith: (self latestComponent).
\end{code}

Se agrega un mensaje que asignará a la variable de instancia una instancia de NullComponent cuando
el valor que devuelva la colaboración del lado derecho sea nil. Esto debe ser así cuando no tenemos
control sobre esas colaboraciones o queremos encapsular la utilización del NullComponent a cierta parte
del código, si se reemplazan todos los lugares donde se utiliza un nil, en lugar de un Component,
por un NullComponent entonces el mensaje de asignación no debería ser necesario.

Luego de realizar las operaciones anteriores la variable de instancia \lstinline{component} no
debería contener en ningún caso nil, solo instancias de Component o NullComponent. El siguiente y
último paso es reemplazar los condicionales por el envío de mensajes polimórficos y el traslado
del comportamiento de nil al NullComponent:

\begin{code}
Context>>launchRocket
    rocketHeight := component ifNil: [0] ifNotNil: [component ignite].

"Luego de aplicar el refactoring"

NullComponent>>ignite
    ^0

Context>>launchRocket
    rocketHeight := component ignite.
\end{code}







\section{Descripción completa de la funcionalidad del refactoring implementado}

El modelo de IntroduceNullObject es más complejo que el de ExtractToMethodObject. Ambos fueron
desarrollados
% REF a TDD
utilizando TDD y la principal diferencia que se encontró es que al empezar a desarrollar de manera
incremental las precondiciones del refactoring eran complejas y requerían que las estructuras
recibidas por el objeto central del modelo, el encargado de aplicar el refactoring, una instancia de
la clase \lstinline{IntroduceNullObject}, fueran también complejas (colleciones con varios niveles
de profundidad). Resultó conveniente repartir las responsabilidades y encapsular las distintas
estructuras que necesita el refactoring en otros objetos que luego éste recibirá. Las precondiciones
chequeadas por el refactoring se redujeron porque los objetos que recibe son válidos entonces las
validaciones espcíficas de esas estructuras están encapsuladas en esos objetos.

El refactoring se utiliza de la siguiente manera:

\begin{code}
refactoring := IntroduceNullObject
    on: 'ivar1'
    of: aRocket class
    usingAsConcreteClass: ConcreteFuelTank
    usingAsNullClass: NullFuelTank
    for: replacementParameters
refactoring apply.
\end{code}

\lstinline{IntroduceNullObject} es subclase de \lstinline{Refactoring} y por lo tanto sobrescribe el
mensaje \lstinline{apply}. Los colaboradores que recibe en el mensaje de creación de instancia son 
los que necesita para poder aplicar el refactoring.

\subsection*{La variable de instancia y la clase}
En el keyword \lstinline{on:} recibe el nombre de la variable de instancia que se reemplazará por el
Null Object. Hasta ahora el refactoring Introduce Null Object parecía aplicarse a un tipo. El tipo
que utilizamos de referencia era Componente y los cambios, reemplazo de asignaciones de nil o 
chequeos por nil, se realizaban sin un scope determinado, es decir, se aplicaban a lo largo de todo
el sistema. Realizar algo así en un sistema dinámico como Cuis no es posible sin sacrificar
demasiada confianza sobre la preservación del comportamiento del refactoring. El proceso para encontrar
todos los lugares en los que se utiliza una instancia de la clase Component no es claro y, aunque lo
fuera, como los tipos son dinámicos la variable podría contener un conjunto de tipos durante su
vida, no solo Component, dependiendo el momento de ejecución. Por último siempre queda el problema de
la metaprogramación que vuelve inútil cualquier intento de inferencia de tipos estática. Por esta
razón decidimos darle al cambio un scope determinado y que el usuario guíe al refactoring con su
conocimiento del modelo.

El scope elegido es el de una clase. La clase provee un equilibrio entre el método, que sería
demasiado restrictivo, probablemente existan pocos casos en los que se justifique introducir este
refactoring solo para eliminar los condicionales de un solo método, y todo el sistema, ya que no
existe un scope claro más grande que la clase o por lo menos Cuis no lo provee ya que las categorías
o los paquetes no permiten ese tipo de funcionalidad. La clase que se utilizará como contexto es la 
recibida en el keyword \lstinline{of:}.

Dentro de una clase se eligirá un colaborador interno o variable de instancia en particular sobre la
cual se realicen chequeos por nil para aplicar el refactoring. Al no contar con un tipo específico
para buscar todas las variables es necesario fijar la o las variables que el programador, por su
conocimiento del modelo, ya sabe que deberían pertenecer a ese tipo. Esto se apoya en las ventajas
ya expuestas de floss refactoring, el caso en el que el programador tiene claro el cambio que quiere
realizar y la herramienta lo automatiza para reducir la posibilidad de errores y aumentar la
velocidad de ejecución. La herramienta no asume entender el modelo ni poder cambiarlo sin guía del
usuario. Se elige como variable a fijar una variable de instancia porque son las variables que están
disponibles a través de todo el scope de la clase y son más comunes que las variables de clase o
las variables de pool.

El refactoring tendrá como objeto central a la variable de instancia elegida, todo los condicionales
modificados serán sobre esa variable. La única precondición que incluye a la variable de
instancia o a la clase es la que verifica que la variable pertenezca a la clase.

\subsection*{La jerarquía de Null Object}

La jerarquía que se asumirá será aquella en la que Component y NullComponent heredan de
AbstractComponent, veremos por qué más adelante. Sin embargo, no es necesario que efectivamente la
implementen, no es una precondición. La única precondición con respecto a las clases Component y
NullComponent es que no sean iguales. Si se utilizara la misma clase no se podrían enviar mensajes
polimórficos que se respondan de forma distinta dependiendo el tipo de la variable.

\subsection*{Los argumentos para los condicionales a reemplazar}

Los condicionales a reemplazar ya están determinados por la variable de instancia y su clase. Smalltalk no
cuenta con condicionales como parte de su sintaxis, éstos se implementan con polimorfismo. Por lo tanto,
los tipos de condicionales soportados por el refactoring no son más que un conjunto mensajes:

\begin{itemize}
    \item Nil Selectors
    \begin{itemize}
        \item \lstinline{ifNil:}
        \item \lstinline{ifNil:ifNotNil:}
        \item \lstinline{ifNotNil:}
        \item \lstinline{ifNotNil:ifNil:}
    \end{itemize}

    \item Boolean Selectors
    \begin{itemize}
        \item \lstinline{ifTrue:}
        \item \lstinline{ifTrue:ifFalse:}
        \item \lstinline{ifFalse:}
        \item \lstinline{ifFalse:ifTrue:}
    \end{itemize}
\end{itemize}

Los Nil Selectors siempre tienen que tener como receptor del mensaje a la variable de instancia y los
Boolean Selectors tienen que tener como receptor del mensaje a un envío de mensaje \lstinline{isNil} a la
variable de instancia:

\begin{code}
"Nil Selector: ejemplo de receptor"
component ifNil: [^4]    

"Boolean Selector: ejemplo de receptor"
component isNil ifTrue: [^4]    
\end{code}

Todos los condicionales soportados serán reemplazados por el envío de un mensaje polimórfico:

\begin{code}
Context>>fuel
    component ifNil: [unit zero] ifNotNil: [component fuelLeftIn: unit]

"Pasará a ser:"
Context>>fuel
    component polymorphicSelector: unit

"Agregando:"
Component>>polymorphicSelector: aUnit
    ^self fuelLeftIn: aUnit

NullComponent>>polymorphicSelector: aUnit
    ^aUnit zero
\end{code}

El usuario deberá elegir el selector polimórfico y los nombres de los parámetros para las variables
a parametrizar. En el ejemplo anterior el selector polimórfico es \lstinline{polymorphicSelector} y
el nombre del parámetro \lstinline{aUnit}.

% DIAGRAM un diagrama de objetos mostrando cómo intereactúan INO y RINWPP

\section{Funcionalidad de ReplaceIfNilWithPolymorphismParameters}

\lstinline{ReplaceIfNilWithPolymorphismParameters} es una clase que representa los argumentos necesarios
para reemplazar un chequeo por nil con polimorfirsmo. Los colaboradores internos de la clase son el
\lstinline{methodNode} del método, la raíz del AST, en el cual se encuentra el chequeo por nil,
el \lstinline{messageNode} correspondiente al chequeo por nil y una instancia de
\lstinline{VariablesToParametrizeKeywordsDefinitions}. Los dos primeros identifican y referencian al
chequeo por nil que reemplazaremos. El tercero es un objeto que contiene la definición del 
mensaje polimórfico: los keywords, qué variable a parametrizar irá en qué keyword y los nombres de
los parámetros.

\lstinline{ReplaceIfNilWithPolymorphismParameters} valida que el \lstinline{messageNode} sea un mensaje
soportado por el refactoring.

% WRITE Explicar cómo se valida que el mensaje sea soportado y las complicaciones que trae eso.
% esto no lo voy a escribir hasta que no haga los tests para ver si ya funciona todo bien o no 
Verificar que el \lstinline{MessageNode} represente un envío de mensaje soportado es un poco más 
complicado que verificar que el selector del nodo pertenezca al conjunto de selectores soportados
y que el receptor del mensaje sea una variable de instancia de la clase del método. La causa son algunas
optimizaciones que se realizan al construir el AST, en específico los MessageNode de un conjunto de mensajes,
dentro de los cuales encontramos a ifNil:ifNotNil: ifNotNil: y ifNotNil:ifNil:.


Además, verifica que la definición del mensaje polimórfico sea consistente con
las variables a parametrizar en los bloques y que no existan colisiones entre
los nombres de los parámetros del mensaje polimórfico y las variables temporales
de los bloques que recibe el chequeo por nil como argumento.


\subsection{Validaciones de la colección de ReplaceIfNilWithPolymorphismParameters}

\lstinline{IntroduceNullObject} valida que los selectores seleccionados para los mensajes polimórficos
no colisionen con otro mensaje de la jerarquía del Null Object. 


\subsection{Aplicación del refactoring}

Al contar con la colección de \lstinline{ReplaceIfNilWithPolymorphismParameters} la aplicación del
refactoring reemplaza cada uno de esos chequeos de nil con lo argumentos correspondientes. Los
cambios que se deben realizar son:

\begin{itemize}
    \item agregar el mensaje polimórfico al Component
    \item agregar el mensaje polimórfico al NullComponent
    \item reemplazar el chequeo de nil por el envío del mensaje polimórfico
\end{itemize}

El proceso para compilar el nuevo método en las clases Component y NullComponent es igual, lo que 
cambia es el bloque utilizado como fuente:

\begin{code}
component ifNil: ['The NullComponent block'] ifNotNil: ['The Component block']
\end{code}

Los mensajes que tienen un solo keyword no cuentan con un bloque para alguna de las dos clases, en
ese caso el método que se compilará será vacío ya que no se requiere ningún comportamiento si el
objeto tiene ese tipo. Veamos un ejemplo:

\begin{code}
Context>>fuel
    ^component ifNotNil: [
        self fuelLevelChecked.
        component fuelLevel]

Component>>fuelLevelWithLog: context
    context fuelLevelChecked.
    ^self fuelLevel

NullComponent>>fuelLevelWithLog: context
    "Empty method"

Context>>fuel
    component fuelLevelWithLog: self
\end{code}

El chequeo por nil se realizó con el mensaje \lstinline{ifNotNil:}, por lo tanto si el objeto es nil
no se espera comportamiento alguno. Por esa razón el método en la clase \lstinline{NullComponent} resulta
vacío.

La generación del código de los métodos para cada caso se realiza reemplazando todas las variables
a parametrizar, todas las variables que se referencian en el bloque excepto las que sean temporales
% WRITE qué sucedía con los argumentos de los bloques? puede ser que no tenían?
del bloque, por el nombre del parámetros correspondiente. Para realizar esto necesitamos obtener el 
código fuente del bloque y luego obtener el rango en el código fuente del método de cada variable a
parametrizar y restarle el offset del inicio del bloque.

Finalmente, se reemplaza el chequeo de nil por el envío del mensaje polimórfico. Se toma el código
fuente del método, se realiza el reemplazo y luego se vuelve a compilar en la clase Context. El
único detalle de este proceso es que la obtención del rango que comprende el chequeo de nil se debe
realizar manualmente porque el \lstinline{Encoder} no posee esta información. En el caso de
\lstinline{MessageNode}s se cuenta con el rango de los keywords y sus argumentos pero no del
receptor. Por lo tanto el rango de un \lstinline{MessageNode} incluyendo al receptor debe tomar el
lado derecho del rango del envío de mensaje y el lado izquierdo del receptor. El receptor en nuestro
caso puede ser una variable o un envío de mensaje (\lstinline{instVar isNil}) por lo cual es simple
obtener el rango consultando al Encoder. La construcción del envío del mensaje la realiza el
\lstinline{ReplaceIfNilWithPolymorphismParameters} correspondiente, ya que cuenta con el mapeo entre
variable a parametrizar y keyword del selector polimórfico. 


\subsection{¿Cómo se utiliza?}

El usuario puede iniciar el refactoring desde el listado de clases en el System Browser porque lo que
necesita el \lstinline{RefactoringApplier} para iniciar la obtención de los argumentos es la clase
a utilizar como contexto:

% DIAGRAM mostrar cómo se selecciona el refactoring

El \lstinline{RefactoringApplier} es el encargado de recabar los argumentos para el
\lstinline{IntroduceNullObject}, crear el refactoring y aplicarlo. Colaborará con el
\lstinline{IntroduceNullObjectForm} y el \lstinline{IntroduceNullObjectRequest} para conseguir esto.
La distribución de responsabilidades quedará reflejada en el siguiente diagrama de secuencia
que muestra cómo colaborane en un caso exitoso:

% DIAGRAM colaboración en caso exitoso

y la experiencia para el usuario sería:

% DIAGRAM mostrar el paso a paso de aplicación del refactoring.

\subsection{Comentarios adicionales}
% WRITE ¿Por qué lo aplicamos solo en el contexto de una clase? ¿Qué otras alternativas se nos ocurrieron
% en un sistema dinámico como Smalltalk en el que no contamos con el tipo o la clase de una variable?
Ya veremos que es fácil pasar de realizar el cambio en el contexto de una clase a realizarlo en todo
el sistema. Bueno quizás no es tan fácil en un sistema dinámico como Smalltalk. Pero tiene ventajas
y desventajas realizarlo de esta manera, si lo vamos aplicando progresivamente podemos identificar
más fácil si algo se rompe y dónde.


% WRITE Precursor del refactoring en referencia al paper de Methodology

\subsection{Preservación del comportamiento}

Aplicaré los refactorings a distintas partes de Cuis y después le voy a correr los tests.


\chapter{Conclusiones}

El trabajo que se realiza también contribuye a la identificaciónd de smells porque tiene que
encontrar todos los casos en los que se produce una asignación a null



%%%% BIBLIOGRAFIA
\backmatter
%\bibliography{tesis}

\end{document}
